% =============================================================================
% Machine Learning & Deep Learning -- Project Proposal Template
% Alpen-Adria-Universität Klagenfurt
%
% Page limit: 4 pages, EXCLUDING references. Everything else counts,
% including figures, tables, the appendix and the AI disclosure statement.
% Do not change margins, font sizes or spacing.
%
% The grey boxes describe what each section must contain. Write your text
% below each box. Before submitting, hide all boxes by changing
% \guidancetrue to \guidancefalse below, then submit the PDF in Moodle.
% =============================================================================

\documentclass{article}

% "preprint" shows your names; without it the template produces an
% anonymous submission.
\usepackage[preprint]{neurips_2025}

\usepackage[utf8]{inputenc} % allow utf-8 input
\usepackage[T1]{fontenc}    % use 8-bit T1 fonts
\usepackage{hyperref}       % hyperlinks
\usepackage{url}            % simple URL typesetting
\usepackage{booktabs}       % professional-quality tables
\usepackage{amsmath}        % mathematical environments
\usepackage{amsfonts}       % blackboard math symbols
\usepackage{nicefrac}       % compact symbols for 1/2, etc.
\usepackage{microtype}      % microtypography
\usepackage{xcolor}         % colors
\usepackage{graphicx}       % figures

% --- Guidance boxes -----------------------------------------------------------
\newif\ifguidance
\guidancetrue % change to \guidancefalse before submission
\newcommand{\guide}[1]{%
  \ifguidance
    \par\noindent\fcolorbox{gray!60}{gray!8}{%
      \parbox{\dimexpr\linewidth-2\fboxsep-2\fboxrule\relax}{\small\color{black!75}#1}}%
    \par\medskip
  \fi}
% -----------------------------------------------------------------------------

\title{Project Proposal: Your Project Title}

\author{%
  First Student \\
  Student ID: 00000000 \\
  \texttt{first.student@edu.aau.at} \\
  \And
  Second Student \\
  Student ID: 00000000 \\
  \texttt{second.student@edu.aau.at} \\
  \And
  Third Student \\
  Student ID: 00000000 \\
  \texttt{third.student@edu.aau.at} \\
}

\begin{document}

\maketitle

\guide{Projects are group work for 2--3 students. The group shares one domain, one dataset and one learning task, and each member takes a significantly different deep learning approach. State clearly what each member will do, e.g.\ ``In this project, (name) will focus on \ldots\ whereas (name) will \ldots''. All sections below are required and must answer all of their questions.}

% =============================================================================
\section{Motivation}

\guide{Why are you doing this project, e.g.\ personal motivation or importance of the problem? What is the goal of your application, and why is it important? What is your hypothesis, i.e.\ why do you believe the problem is solvable?}

Your text.

% =============================================================================
\section{Learning Task}

\subsection{Training Experience}

\guide{Where does your training data come from? Why do you think it can actually be used to solve the problem? Show examples of your data and use them to explain the learning task.}

Your text.

\subsection{Task}

\guide{What are you going to do? Which type of learning suits the task?}

Your text.

\subsection{Performance Measure}

\guide{How will you measure success? Which measures will you use, and why? How do you interpret their values for your problem? For example, is a precision of 0.7 good or bad?}

Your text.

% =============================================================================
\section{Previous Work}

\guide{Summarize two to three papers, blog posts or articles that consider a learning task (problem, dataset, measures) similar to yours. How can you apply their approach to your problem? What modifications might be required? Can you use it as a baseline in your evaluation? Cite your sources, e.g.\ \citet{bishop2024deep} or \citep{goodfellow2016deep}, and add them to \texttt{references.bib}.}

Your text.

% =============================================================================
\section{Plan}

\guide{How will you solve the learning task with (an extension of) state-of-the-art methods? Describe the approaches in enough detail to show their benefits, how they differ from previous work, and that they can be implemented within the project. State which approach each group member takes, e.g.\ in a table like Table~\ref{tab:roles}.}

Your text.

\begin{table}[h]
  \caption{Approaches of the group members.}
  \label{tab:roles}
  \centering
  \begin{tabular}{ll}
    \toprule
    Group member & Approach \\
    \midrule
    First Student  & e.g.\ convolutional network trained from scratch \\
    Second Student & e.g.\ pre-trained network fine-tuned to the task \\
    Third Student  & e.g.\ vision transformer \\
    \bottomrule
  \end{tabular}
\end{table}

% =============================================================================
\section{Risk Management}

\guide{What is your ``plan B'' if something does not go as expected, e.g.\ no data, poor-quality data, or training that does not converge?}

Your text.

% =============================================================================
\bibliographystyle{plainnat}
\bibliography{references}

% =============================================================================
\appendix

\section{AI Disclosure Statement}

\guide{Declare any use of generative AI tools and coding assistants. State which tools you used and for what purpose, e.g.\ brainstorming, literature search, code generation or text drafting. If you did not use any, say so.}

Your text.

\end{document}
